\markdownRendererDocumentBegin
\markdownRendererWarning{The `hybrid` option has been soft-deprecated.}{The `hybrid` option has been soft-deprecated.}{Consider using one of the following better options for mixing TeX and markdown: `contentBlocks`, `rawAttribute`, `texComments`, `texMathDollars`, `texMathSingleBackslash`, and `texMathDoubleBackslash`. For more information, see the user manual at <https://witiko.github.io/markdown/>.}{Consider using one of the following better options for mixing TeX and markdown: `contentBlocks`, `rawAttribute`, `texComments`, `texMathDollars`, `texMathSingleBackslash`, and `texMathDoubleBackslash`. For more information, see the user manual at <https://witiko.github.io/markdown/>.}\markdownRendererSectionBegin
\markdownRendererHeadingOne{Exponenciación Binaria}\markdownRendererInterblockSeparator
{}La operación $a^b$ en la gran mayoría de librerías se implementa en $O(b)$.\markdownRendererHardLineBreak
{}Sin embargo, si representamos $b$ en base 2 (binario), se puede realizar en $O(log b)$:\markdownRendererInterblockSeparator
{}\markdownRendererInputFencedCode{./_markdown_main/745e27c31ab3cb2b79f8ec72c1a6ff98.verbatim}{cpp}{cpp}\markdownRendererInterblockSeparator
{}
\markdownRendererSectionEnd \markdownRendererSectionBegin
\markdownRendererHeadingOne{Aritmética Modular}\markdownRendererInterblockSeparator
{}Cuando los números de entrada son muy grandes, los problemas suelen requerir que apliquemos aritmética modular con un número p, típicamente $10^9$ + 7 (1000000007).\markdownRendererParagraphSeparator
{}Por ejemplo, supongamos que queremos calcular: $r = (a * b * c)\markdownRendererParagraphSeparator
{}Si realizamos la multiplicación directamente, los números intermedios pueden exceder la capacidad de almacenamiento, aunque matemáticamente sea correcto. Por eso usamos operaciones modulares paso a paso.\markdownRendererInterblockSeparator
{}\markdownRendererSectionBegin
\markdownRendererSectionBegin
\markdownRendererHeadingThree{Operaciones mediante aritmética modular}\markdownRendererParagraphSeparator
{}\markdownRendererInputFencedCode{./_markdown_main/d9bb1a7f3f6d6bf7b5df733d7aafc9b4.verbatim}{cpp}{cpp}
\markdownRendererSectionEnd 
\markdownRendererSectionEnd 
\markdownRendererSectionEnd \markdownRendererDocumentEnd